\section{Pendahuluan}\label{sec:intro}
\subsection{Pertanyaan ekonomi}
Depresiasi sering dihubungkan dengan peningkatan daya saing karena harga barang domestik menjadi relatif lebih murah bagi pembeli asing. Penjelasan tersebut merupakan titik awal yang berguna, tetapi belum menetapkan arah perubahan neraca perdagangan. Nilai ekspor dan impor juga bergantung pada harga internasional, permintaan, kemampuan produksi, dan waktu penyesuaian. Respons suatu barang atau pasar dapat berbeda dari respons agregat. Penelitian ini menanyakan sejauh mana perbaikan yang diharapkan terlihat pada perdagangan Indonesia, dan apakah kesimpulannya bertahan ketika arah perubahan kurs serta komposisi barang dipisahkan.

Motivasi awal berkaitan dengan besarnya peran komoditas dalam ekspor Indonesia. Pada rincian HS-6 yang digunakan, komoditas yang dapat dipetakan secara konservatif mencakup \CommodityShare\% nilai ekspor 2025, sedangkan manufaktur nonkomoditas \ManufacturingShare\%. Penetapan harga yang mengikuti pasar internasional memberi alasan untuk menilai apakah perubahan kurs domestik mempunyai hubungan yang sama pada kedua kelompok. Namun komposisi nilai tidak langsung memberikan ukuran sensitivitas kurs. Pangsa yang besar dan koefisien yang kecil menjawab pertanyaan berbeda, sehingga bukti deskriptif harus dibaca bersama estimasi dan ketidakpastiannya.

Penelitian menggunakan neraca relatif sebagai outcome, mengikuti paper acuan. Implikasinya perlu dipertahankan sepanjang pembahasan: perbaikan rasio dapat terjadi ketika ekspor naik, ketika impor turun, atau ketika keduanya bergerak dengan kecepatan berbeda. Model yang ada belum menetapkan bahwa perubahan berasal khusus dari ekspor. Hipotesis mengenai kandungan bahan baku impor dalam ekspor sengaja tidak dimasukkan dalam run ini.

\subsection{Cakupan dan pertanyaan yang dapat dijawab}
Lapis pertama adalah ekspor Indonesia ke seluruh negara dan impor Indonesia dari seluruh negara. Lapis kedua memisahkan arus dengan Australia, Tiongkok, India, Jepang, Korea Selatan, Malaysia, Singapura, Thailand, Amerika Serikat, dan Vietnam. Lapis ketiga membentuk tujuh kelompok barang untuk agregat seluruh mitra maupun masing-masing bilateral. Istilah agregat tidak merujuk pada penjumlahan perdagangan antarsemua negara di dunia.

Data tersedia pada HS-6, tetapi unit regresi barang adalah kelompok hasil klasifikasi. Terdapat 77 seri kelompok: tujuh kelompok pada sebelas cakupan pasar. Masing-masing diestimasi menggunakan ARDL dan NARDL. Bersama satu seri agregat dan sepuluh bilateral, struktur itu menghasilkan 88 seri utama dan 176 estimasi. Jumlah spesifikasi tersebut bukan jumlah pengamatan independen. Kelompok bertumpang tindih dan memakai banyak observasi yang sama.

Tiga pertanyaan dibedakan. Pertama, apakah hubungan depresiasi dengan rasio neraca mengarah pada perbaikan atau pemburukan? Kedua, apakah terlihat pemburukan awal sebelum perbaikan berikutnya? Ketiga, apakah model dan interval mendukung penafsiran pola itu? Pemisahan ini mencegah bentuk grafik, signifikansi satu koefisien, dan kelayakan hubungan level diperlakukan sebagai bukti yang saling menggantikan.

\subsection{Literatur acuan dan posisi analisis}
Bahmani-Oskooee dan Fariditavana \cite{bahmani} menilai J-curve melalui hubungan jangka pendek dan jangka panjang serta membandingkan penyesuaian linier dengan penyesuaian yang memisahkan apresiasi dan depresiasi. Paper pengguna memakai enam mitra Amerika Serikat pada data kuartalan 1971Q1--2013Q3. Kerangka tersebut mendasari penggunaan rasio impor/ekspor, aktivitas domestik dan mitra, serta kurs riil. Penerapan Indonesia memperluas disaggregasi ke kelompok barang pada sampel yang lebih pendek.

Pesaran, Shin, dan Smith \cite{pss} menyediakan pendekatan bounds untuk menilai hubungan level ketika integrasi regressor tidak harus seragam antara orde nol dan satu. Shin, Yu, dan Greenwood-Nimmo \cite{shin} mengembangkan pemisahan perubahan positif dan negatif serta multiplier dinamis nonlinear. Dalam penelitian ini pemisahan digunakan untuk menilai apakah akumulasi depresiasi dan apresiasi mempunyai hubungan berbeda. NARDL memungkinkan asimetri; ia tidak mengharuskan asimetri ditemukan atau membuat kurs menjadi eksogen.

Implementasi menggunakan paket Stata dan prosedur yang dijelaskan Kripfganz dan Schneider \cite{ks}. Lag, representasi error correction, dan bounds finite-sample diperiksa pada tiap kasus. Panjang lag serta beberapa aturan penyaringan disesuaikan dengan sampel Indonesia. Istilah yang tepat adalah adaptasi metode paper, bukan reproduksi identik seluruh keputusan empirisnya.

Nilai tambah analisis terletak pada konsistensi definisi, perbandingan arah kurs, dan pelaporan lengkap kasus yang lebih maupun kurang mendukung. Data belum menyediakan identifikasi kausal ataupun pengukuran volume untuk memisahkan mekanisme harga dan kuantitas. Penafsiran ekonomi karena itu harus dibedakan dari mekanisme yang telah diidentifikasi.
